\title{DeepSeekMath: Pushing the Limits of Mathematical Reasoning in Open Language Models}

\begin{abstract}
The structured and intricate nature of mathematical reasoning creates considerable difficulties for language models. This study presents DeepSeekMath~7B, developed by further pre-training DeepSeek-Coder-Base-v1.5 7B on a dataset consisting of 120B tokens pertaining to mathematics, extracted from Common Crawl alongside both natural language and programming code corpora. DeepSeekMath~7B attains an impressive 51.7\% on the advanced MATH benchmark, demonstrating strong capability without dependence on ensemble methods or external computational tools, and closely rivaling the results of Gemini-Ultra and GPT-4. By sampling 64 solutions for self-consistency, DeepSeekMath~7B achieves a performance of 60.9\% on MATH. This enhanced proficiency in mathematical reasoning arises from two principal factors: (1) the careful exploitation of large-scale, publicly accessible web resources enabled by a rigorously designed data selection process, and (2) the proposal of Group Relative Policy Optimization (GRPO), an adaptation of Proximal Policy Optimization (PPO), which augments mathematical reasoning ability while also improving the efficiency of PPO in terms of memory consumption.
\end{abstract}

\section{Introduction}

Large language models (LLMs) have transformed methods for mathematical reasoning in artificial intelligence, leading to notable progress on benchmarks in both quantitative reasoning \citep{MATH} and geometric problem solving \citep{trinh2024solving}. These models have also demonstrated considerable value in supporting human efforts to tackle challenging mathematical tasks \citep{tao}. Despite these advancements, the most advanced systems—such as GPT-4 \citep{gpt4} and Gemini-Ultra \citep{gemini}—remain inaccessible to the public, while available open-source alternatives still fall significantly short in performance.

This paper presents DeepSeekMath, a specialized language model tailored for mathematical tasks, which considerably surpasses the mathematical proficiency of existing open-source models and nearly matches GPT-4’s results on academic evaluation sets. Central to this achievement is the construction of the DeepSeekMath Corpus, an extensive and meticulously curated pre-training resource containing 120B mathematical tokens. This corpus is assembled from Common Crawl (CC) data using a fastText-based classifier \citep{joulin2016fasttext}. The classifier is initially trained on positive samples drawn from OpenWebMath \citep{openwebmath}, with a varied assortment of web pages designated as negative instances. Following this initial phase, the classifier is used to identify further relevant data within CC, which is subsequently fine-tuned through manual annotation. This enhanced set serves to further retrain the classifier, boosting its extraction accuracy. Empirical results show that the generated corpus is of high caliber: our foundational DeepSeekMath-Base 7B model attains 64.2\% accuracy on GSM8K \citep{gsm8k} and 36.2\% on the advanced MATH benchmark \citep{MATH}, achieving better performance than Minerva 540B \citep{minerva}. Furthermore, because the DeepSeekMath Corpus includes multilingual data, we also observe gains on Chinese mathematical test sets \citep{wei2023cmath,agieval}. We propose that our efforts in constructing mathematical datasets mark an initial step for further exploration and refinement within the research community.

DeepSeekMath-Base is built upon DeepSeek-Coder-Base-v1.5 7B \citep{deepseek-coder}, as our findings suggest that initiating training from a code-centric model yields superior results relative to employing a general LLM foundation. In addition, our experiments reveal that the process of math-focused training elevates the model's performance on MMLU \citep{mmlu} and BBH benchmarks \citep{bbh}, demonstrating that such training not only strengthens mathematical proficiency but also broadens the model’s overall reasoning skills.

Following the pre-training phase, we perform mathematical instruction tuning on DeepSeekMath-Base utilizing datasets focused on chain-of-thought \citep{cot}, program-of-thought \citep{pot,pal}, and tool-integrated reasoning \citep{tora}. The tuned model, DeepSeekMath-Instruct 7B, surpasses all other 7B models and achieves performance on par with 70B-scale open-source models that have undergone instruction tuning.

In addition, we present Group Relative Policy Optimization (GRPO), a novel reinforcement learning (RL) algorithm derived from Proximal Policy Optimization (PPO) \citep{schulman2017proximal}. Unlike PPO, GRPO dispenses with the use of a critic model, instead utilizing group scores to estimate the baseline, thereby leading to significant reductions in required training resources. Leveraging only a subset of English instruction tuning data, GRPO achieves marked gains over the robust DeepSeekMath-Instruct model, both for in-domain datasets (GSM8K: 82.9\% $\rightarrow$ 88.2\%, MATH: 46.8\% $\rightarrow$ 51.7\%) and for out-of-domain scenarios, such as with CMATH (84.6\% $\rightarrow$ 88.8\%) during RL training. We additionally offer a coherent framework to interpret and compare various techniques—including Rejection Sampling Fine-Tuning (RFT) \citep{yuan2023scaling}, Direct Preference Optimization (DPO) \citep{dpo}, PPO, and GRPO—showing within this unified perspective that these methods can essentially be viewed as forms of direct or streamlined RL. To investigate key components within this paradigm, we conduct comprehensive experiments examining factors such as online versus offline training, process versus outcome supervision, and comparisons between single-turn and iterative RL approaches. Finally, we elucidate the mechanisms by which our RL methods enhance the capabilities of instruction-tuned models and outline avenues for future advances in effective RL, grounded in this unified framework.

\subsection{Contributions}
We present scalable methods for math pre-training and conduct a comprehensive investigation and evaluation of reinforcement learning techniques.

\textbf{Math Pre-Training at Scale}

\begin{itemize}[topsep=0pt]
    \item This study demonstrates that the publicly available Common Crawl dataset harbors significant resources pertinent to mathematical applications. Through the application of a carefully crafted data selection pipeline, we assembled the DeepSeekMath Corpus—a high-quality collection comprising 120B tokens sourced from web pages curated for mathematical relevance. This corpus is nearly seven times larger than the math web page subset employed by Minerva \citep{minerva} and nine times larger than the recently introduced OpenWebMath \citep{openwebmath}.

\item DeepSeekMath-Base 7B, our pre-trained foundational model, matches the performance level of Minerva 540B \citep{minerva}, demonstrating that parameter count alone does not solely determine mathematical reasoning proficiency. This result suggests that even models with fewer parameters can exhibit robust capabilities when trained on high-quality datasets.

\item We present results from experiments focused on mathematical training. Pre-training models on code before exposing them to mathematical content enhances their performance in solving math problems, whether tools are employed or not. This provides some insight into the ongoing debate over whether code training enhances reasoning capabilities; our findings suggest that, in the context of mathematical reasoning, it indeed does.

\item While the use of arXiv papers as training data is widespread—particularly within the context of mathematical research—this approach yields no significant gains across the mathematical benchmarks evaluated in this study.
\end{itemize}

\textbf{Exploration and Analysis of Reinforcement Learning}

\begin{itemize}[topsep=0pt]
    \item We present Group Relative Policy Optimization (GRPO), a novel reinforcement learning approach that bypasses the use of a critic model. GRPO utilizes group scores to estimate baselines, which leads to a substantial decrease in required training resources in comparison to Proximal Policy Optimization (PPO).
    \item Our results show that GRPO yields substantial improvements in the instruction-tuned model DeepSeekMath-Instruct, utilizing only instruction-tuning data. In addition, we note improved performance on out-of-domain tasks throughout reinforcement learning.
    \item We offer a cohesive framework for understanding various approaches, including RFT, DPO, PPO, and GRPO. Through comprehensive experimentation—such as contrasting online versus offline training, evaluating outcome versus process supervision, and comparing single-turn with iterative reinforcement learning—we conduct an in-depth analysis of the crucial factors underlying this framework.
    \item Drawing on this unified framework, we investigate the mechanisms that make reinforcement learning effective, and propose several future research avenues for enhancing the efficiency of reinforcement learning with LLMs.
\end{itemize}

\subsection{Summary of Evaluations and Metrics}

\begin{itemize}[topsep=0pt]
    \item \textbf{English and Chinese Mathematical Reasoning}: We carry out extensive evaluations of our models on both English and Chinese mathematical benchmarks, encompassing questions ranging from elementary to collegiate difficulty. For English, we utilize datasets such as GSM8K \citep{gsm8k}, MATH \citep{MATH}, SAT \citep{llemma}, OCW Courses \citep{minerva}, and MMLU-STEM \citep{mmlu}. The Chinese benchmarks employed are MGSM-zh \citep{mgsm}, CMATH \citep{wei2023cmath}, Gaokao-MathCloze \citep{agieval}, and Gaokao-MathQA \citep{agieval}. We assess the models in two aspects: the capacity to independently produce comprehensive, self-explanatory solutions and the competence in resolving problems with the aid of Python.

On English benchmark evaluations, DeepSeekMath-Base demonstrates performance on par with the proprietary Minerva 540B model \citep{minerva}, and consistently outperforms all available open-source base models—such as Mistral 7B \citep{mistral} and Llemma-34B \citep{llemma}—regardless of their exposure to math-specific pre-training, and frequently does so by a substantial margin. Importantly, DeepSeekMath-Base achieves even greater success on Chinese language benchmarks. This advantage likely arises because, unlike earlier studies \citep{minerva, llemma} that exclusively utilize English-language data for math pre-training, our approach also incorporates high-quality data in languages other than English. Subsequent instruction fine-tuning and the application of reinforcement learning yields DeepSeekMath-Instruct and DeepSeekMath-RL, both of which achieve exceptional results, surpassing 50\% accuracy for the first time in open-source models on the advanced, competition-level MATH dataset.

\item \textbf{Formal Mathematics}: DeepSeekMath-Base is assessed on its capacity for informal-to-formal theorem proving via the task described in \citep{dsp_proof}, applied to the miniF2F \citep{minif2f} benchmark, utilizing Isabelle \citep{isabelle} as the designated proof assistant. The results indicate that DeepSeekMath-Base achieves robust few-shot autoformalization results.
\item \textbf{Natural Language Understanding, Reasoning, and Code}: To gain an in-depth view of the model's proficiencies across general understanding, reasoning, and coding, we subject DeepSeekMath-Base to the Massive Multitask Language Understanding (MMLU) benchmark \citep{mmlu}, featuring 57 multiple-choice tasks spanning a broad spectrum of fields, as well as BIG-Bench Hard (BBH) \citep{bbh}, which offers 23 complex tasks demanding multi-step reasoning. In addition, HumanEval \citep{codex} and MBPP \citep{mbpp} serve as primary measures for assessing code generation. Math-oriented pre-training leads to noticeable gains in both reasoning and language comprehension.

\section{Math Pre-Training}

\subsection{Data Collection and Decontamination} This section describes the methodology employed to assemble the DeepSeekMath Corpus utilizing data from Common Crawl. Figure \ref{fig:data_collect} illustrates the iterative workflow developed for methodically curating an extensive mathematics-focused dataset from Common Crawl, commencing with an initial seed corpus composed of a modest yet high-quality set of mathematical data. It should be emphasized that this methodology is also transferable to additional fields, such as coding.

We begin by selecting OpenWebMath \citep{openwebmath}, an assemblage of curated, high-quality mathematical web resources, as our foundational seed dataset. Utilizing this collection, we train a fastText model \citep{joulin2016fasttext} to retrieve additional web pages that are thematically similar to those found in OpenWebMath. For this purpose, 500,000 randomly sampled entries from the seed set serve as positive samples, while an equal number of web pages drawn from Common Crawl are used as negative examples. The training leverages an open-source fastText implementation\footnote{\url{https://fasttext.cc}}, with hyperparameters set as follows: 256 for vector dimensionality, a learning rate of 0.1, a maximum word n-gram length of 3, a minimum word count threshold of 3, and three epochs. To manage the immense scale of the initial Common Crawl dataset, we apply deduplication based both on URLs and on near-duplicate content, yielding a corpus of 40B HTML web pages. Subsequently, the trained fastText model identifies mathematical web documents from the deduplicated set. To ensure the quality of the extracted mathematical materials, we rank retrieved web pages using the scores output by the model, retaining only the top-ranked portion. We then assess how much content to preserve by conducting pre-training experiments on corpora containing the leading 40B, 80B, 120B, and 160B tokens according to model ranking. In this initial phase, we select the top 40B tokens for further use.

Following the initial round of data collection, a significant number of mathematical web pages are still missing from our corpus, primarily due to the limited diversity of positive examples used in training the fastText model. To address this shortcoming, we expand the seed corpus by incorporating new sources of mathematical web content, enhancing the quality of our training data for the fastText model. Our process begins by partitioning all content in Common Crawl into separate domains, where each domain consists of web pages sharing the same root URL. For each of these domains, we compute the proportion of pages that were retrieved during the first data collection iteration. Domains in which more than 10\% of the web pages have already been collected are categorized as math-related domains (for instance, \url{mathoverflow.net}). 

We then proceed to manually label URLs within these domains that correspond to mathematical content (e.g., \url{mathoverflow.net/questions}). Web pages associated with such annotated URLs that have not yet been harvested are subsequently added to the seed corpus. This methodology allows us to incorporate a greater variety of positive examples, ultimately resulting in a better-trained fastText model with improved recall capabilities for mathematical content in the next cycle. After conducting four rounds of data collection, our resulting dataset comprises 35.5M mathematical web pages with a total of 120B tokens. By the fourth iteration, we observe that approximately 98\% of the newly collected data had already been obtained during the third iteration, prompting us to discontinue further collection efforts.

In order to prevent benchmark contamination, we adopt the filtering procedure described in \cite{deepseek-coder}, excluding web content that features questions or answers from notable English mathematical benchmarks such as GSM8K \citep{gsm8k} and MATH \citep{MATH}, as well as Chinese datasets like CMATH \citep{wei2023cmath} and AGIEval \citep{agieval}. Specifically, our approach removes any portion of text within our math training set that contains a 10-gram sequence matching verbatim any substring present in these evaluation benchmarks. For instances where benchmark excerpts are less than 10 grams but at least 3 grams in length, we apply exact string matching to exclude these potentially contaminated web documents.

\subsection{Validating the Quality of the DeepSeekMath~Corpus}

We run pre-training experiments to investigate how the DeepSeekMath~Corpus is compared with the recently released math-training corpora:

\begin{itemize}[topsep=0pt]
    \item \textbf{MathPile} \citep{mathpile}: This dataset compiles 8.9B tokens from varied sources such as textbooks, Wikipedia, ProofWiki, CommonCrawl, StackExchange, and arXiv, with arXiv alone accounting for more than 85\% of the overall content;
    \item \textbf{OpenWebMath} \citep{openwebmath}: Contains 13.6B tokens selected from CommonCrawl, filtered specifically to retain mathematical material;
    \item \textbf{Proof-Pile-2} \citep{llemma}: An aggregate mathematical dataset formed from OpenWebMath, AlgebraicStack (comprising 10.3B tokens of math-related code), as well as arXiv (contributing 28.0B tokens). Our experiments with Proof-Pile-2 adhere to the data proportion of 2:4:1 for arXiv:Web:Code, as specified in \cite{llemma}.
\end{itemize}

\subsubsection{Training Setting}
\label{sec:quality-policy}
We perform math-specific training on a 1.3B-parameter general-purpose language model structured identically to the DeepSeek LLMs \citep{deepseek-llm}, hereafter referred to as DeepSeek-LLM 1.3B. Each mathematical dataset is used to train a distinct model for a total of 150B tokens. All training runs leverage the streamlined and resource-efficient HAI-LLM framework \citep{haillm}. In alignment with the training methodology employed for DeepSeek LLMs, we utilize the AdamW optimizer \citep{adamW}, setting $\beta_1=0.9$, $\beta_2=0.95$, and $\mathrm{weight\_decay}=0.1$. The learning rate scheduling follows a staged approach: an initial ramp-up to its maximum over 2,000 warmup iterations, a reduction to 31.6\% of the peak at the 80\% point in training, and finally a further decay to 10.0\% of the peak by 90\% completion. The maximum learning rate is fixed at $5.3\times10^{-4}$. Training employs a batch size containing 4 million tokens and supports up to 4,000 tokens per input sequence.

\subsubsection{Evaluation Results}

\textbf{The DeepSeekMath~Corpus is of high quality, covers multilingual mathematical content, and is the largest in size.}

\begin{itemize}[topsep=0pt]
    \item \textbf{High-quality}: Downstream results on eight mathematical benchmarks are assessed via few-shot chain-of-thought prompting~\cite{cot}. Table~\ref{tab:corpora_comparison} presents clear evidence that models trained using the DeepSeekMath~Corpus consistently outperform those utilizing alternative corpora. Further, Figure~\ref{fig:corpora_comparisons} illustrates the superior capabilities of the DeepSeekMath-trained model over the Proof-Pile-2-trained model when trained with 50B tokens (corresponding to one epoch of Proof-Pile-2), suggesting that the DeepSeekMath~Corpus is of generally higher quality.
    \item \textbf{Multilingual}: Incorporating multilingual content, the DeepSeekMath~Corpus consists mainly of English and Chinese data, with these two languages forming the core of its dataset. Results in Table~\ref{tab:corpora_comparison} indicate that models trained on DeepSeekMath~Corpus attain improved mathematical reasoning in both English and Chinese. By contrast, mathematical corpora that predominantly focus on English exhibit only minimal gains—or even adverse effects—in reasoning tasks involving Chinese text.
    \item \textbf{Large-scale}: The DeepSeekMath~Corpus greatly surpasses existing mathematical corpora in scale. As can be seen in Figure~\ref{fig:corpora_comparisons}, training DeepSeek-LLM 1.3B with DeepSeekMath~Corpus results in more rapid and sustained learning improvements compared to baselines. Meanwhile, baseline datasets—considerably smaller in size—reach saturation quickly, with training having cycled through the data multiple times.
\end{itemize}

\subsection{Training and Evaluating DeepSeekMath-Base 7B}

This section details DeepSeekMath-Base 7B, a foundational model distinguished by advanced mathematical reasoning skills. The initialization is based on DeepSeek-Coder-Base-v1.5 7B \citep{deepseek-coder}, and training occurs over 500B tokens. The training corpus comprises 56\% DeepSeekMath Corpus, 4\% AlgebraicStack, 10\% arXiv, 20\% code from Github, with the last 10\% drawn from Common Crawl’s English and Chinese natural language content. Apart from using the training configuration outlined in Section \ref{sec:quality-policy}, two key modifications are introduced: the learning rate’s maximum value is adjusted to 4.2e-4, and training employs a batch size of 10M tokens.

In this work, we thoroughly evaluate the mathematical proficiency of DeepSeekMath-Base 7B. Our analysis centers on the model's effectiveness in independently generating complete mathematical solutions, applying external tools to address mathematical tasks, and executing formal theorem proofs. Additionally, we broaden our examination to encompass a general overview of the base model's competencies, such as its aptitude in natural language comprehension, logical reasoning, and programming abilities.

\paragraph{Mathematical Problem Solving with Step-by-Step Reasoning}

We assess the capabilities of DeepSeekMath-Base in mathematical problem-solving with few-shot chain-of-thought prompting \citep{cot}, utilizing eight different benchmarks in both English and Chinese. These datasets include tasks that require quantitative reasoning (such as GSM8K \citep{gsm8k}, MATH \citep{MATH}, and CMATH \citep{wei2023cmath}), as well as multiple-choice formats (including MMLU-STEM \citep{mmlu} and Gaokao-MathQA \citep{agieval}), spanning a broad range of mathematical topics from basic to advanced undergraduate levels.

As detailed in Table \ref{tab:base_math_cot}, DeepSeekMath-Base 7B outperforms all other open-source base models on the full suite of eight benchmarks, including the popular general-purpose Mistral 7B \citep{mistral} and the more recent Llemma 34B \citep{llemma}, the latter of which received mathematical training using Proof-Pile-2 \citep{llemma}. Of particular significance is DeepSeekMath-Base's performance on the MATH dataset, where it exceeds the accuracy of other open-source base models by more than 10 percentage points, and even surpasses Minerva 540B \citep{minerva}—a closed-source model that is 77 times larger, developed atop PaLM \citep{palm} and specifically further trained on mathematical materials.

\paragraph{Mathematical Problem Solving with Tool Use} We assess mathematical reasoning facilitated by program usage on the GSM8K and MATH datasets through few-shot program-of-thought prompting \citep{pot,pal}. To tackle each problem, models generate a Python script, making use of libraries like \textit{math} and \textit{sympy} for advanced calculations. The outcome produced by running the code is treated as the final answer. Table \ref{tab:base_math_pot_proof} demonstrates that DeepSeekMath-Base 7B surpasses the previously leading Llemma 34B in performance.

\paragraph{Formal Mathematics}

Automating formal proofs has proven advantageous for ensuring the correctness and dependability of mathematical reasoning, as well as improving productivity—a trend that has garnered significant scholarly interest recently. In our study, we assess DeepSeekMath-Base 7B using the informal-to-formal proof generation task outlined by \citep{dsp_proof}, where the goal is to construct a formal proof given an informal statement, its formal analog, and a corresponding informal proof. Our experiments employ the miniF2F benchmark \citep{minif2f}, which assesses formal reasoning abilities at Olympiad-level mathematics. For each problem in this suite, we prompt the model with several examples to generate formal proofs in Isabelle. Adhering to the evaluation framework in \cite{dsp_proof}, we utilize the model to create proof outlines and then apply the standard automated prover Sledgehammer \citep{sledgehammer} to complete any unfinished proof components. As indicated in Table \ref{tab:base_math_pot_proof}, DeepSeekMath-Base 7B achieves impressive results in the area of proof autoformalization.

\paragraph{Natural Language Understanding, Reasoning, and Code}

We assess the natural language understanding capabilities using MMLU \citep{mmlu}, reasoning ability via BBH \citep{bbh}, and programming proficiency on HumanEval \citep{codex} and MBPP \citep{mbpp}. According to the results reported in Table \ref{tab:understanding_reasoning_code}, DeepSeekMath-Base 7B demonstrates marked improvements over its predecessor, DeepSeek-Coder-Base-v1.5 \citep{deepseek-coder}, on both MMLU and BBH, underscoring the beneficial role of mathematical training for language understanding and reasoning tasks. Furthermore, by incorporating code-related data during continued training, DeepSeekMath-Base 7B successfully preserves the coding benchmark performance of DeepSeek-Coder-Base-v1.5 on HumanEval and MBPP. In summary, DeepSeekMath-Base 7B notably surpasses the general-purpose Mistral 7B \citep{mistral} across all three benchmarks in reasoning and coding.

\section{Supervised Fine-Tuning}

\subsection{SFT Data Curation}

We develop an instruction-tuning dataset for mathematics that encompasses English and Chinese questions drawn from a range of mathematical domains and spanning multiple levels of difficulty. Each problem is accompanied by solutions presented in chain-of-thought (CoT) \citep{cot}, program-of-thought (PoT) \citep{pot,pal}, and tool-integrated reasoning styles \citep{tora}. Altogether, the dataset consists of 776K training instances.

\begin{itemize}[topsep=0pt]
    \item \textbf{English mathematical datasets}: We provide annotations with tool-assisted solutions for both GSM8K and MATH problems. Additionally, we utilize a subset of MathInstruct \citep{MathInstruct} and include the training portion of Lila-OOD \citep{lila}, comprising problems tackled via CoT or PoT methodologies. The English dataset encompasses a broad spectrum of mathematical domains such as algebra, probability, number theory, calculus, and geometry.
    \item \textbf{Chinese mathematical datasets}: We assemble a collection of Chinese K-12 math problems, covering 76 distinct sub-topics including, for example, linear equations. For each, we supply solutions annotated with both CoT and tool-based reasoning strategies.
\end{itemize}

\subsection{Training and Evaluating DeepSeekMath-Instruct 7B}

Within this section, we present DeepSeekMath-Instruct 7B, a model derived from mathematical instruction tuning using DeepSeekMath-Base as its foundation. During preparation, training instances are randomly joined together until the sequence attains a maximum length of 4K tokens. The model is subsequently trained for 500 steps, utilizing a batch size of 256 and maintaining a fixed learning rate of 5e-5.

We evaluate models' mathematical performance both without and with tool use, on 4 quantitative reasoning benchmarks in English and Chinese. We benchmark our model against the leading models of the time:

\begin{itemize}[topsep=0pt]
    \item \textbf{Closed-source models} comprise the following: (1) the GPT family, with GPT-4 \citep{gpt4} and GPT-4 Code Interpreter~\footnote{\url{https://openai.com/blog/chatgpt-plugins\#\#code-interpreter}} standing out as the most advanced variants, (2) Gemini Ultra and Pro \citep{gemini}, (3) Inflection-2 \citep{inflection-2}, (4) Grok-1~\footnote{\url{https://x.ai/model-card}}, along with a series of models newly launched by Chinese companies, including (5) Baichuan-3~\footnote{\url{https://www.baichuan-ai.com}}, and (6) the state-of-the-art GLM-4~\footnote{\url{https://open.bigmodel.cn/dev/api\#glm-4}}

    from the GLM family \citep{glm}.

These models are designed for broad application and the majority have been subjected to various alignment strategies.

\item \textbf{Open-source models} consist of: (1) DeepSeek-LLM-Chat 67B \citep{deepseek-llm}, (2) Qwen 72B \citep{qwen}, (3) SeaLLM-v2 7B \citep{seallm}, and (4) ChatGLM3 6B \citep{chatglm3}, all of which serve as general-purpose systems. There are also models tailored for mathematical tasks: (5) InternLM2-Math 20B~\footnote{\url{https://github.com/InternLM/InternLM-Math}}, developed from InternLM2 with subsequent fine-tuning on mathematics data followed by instruction-based tuning; (6) Math-Shepherd-Mistral 7B, which adapts the Mistral 7B \citep{mistral} model by employing PPO training \citep{schulman2017proximal} combined with a process supervision reward model; (7) the WizardMath suite \citep{wizardmath}, which advances mathematical capabilities in Mistral 7B and Llama-2 70B \citep{llama2} using evolve-instruct—an approach to instruction tuning utilizing AI-generated instructions—and PPO training, with core data drawn from GSM8K and MATH datasets; (8) MetaMath 70B \citep{metamath}, produced by additional tuning of Llama-2 70B with enhanced versions of the GSM8K and MATH corpora; (9) ToRA 34B \cite{tora}, which involves fine-tuning CodeLlama 34B to perform mathematical reasoning integrated with tools; and (10) MAmmoTH 70B \citep{MathInstruct}, which leverages instruction tuning on Llama-2 70B utilizing MathInstruct data.

As presented in Table \ref{tab:sft_rl_math}, when tool use is restricted during evaluation, DeepSeekMath-Instruct 7B exhibits impressive step-by-step reasoning abilities. Specifically, on the competition-level MATH dataset, our model achieves a lead of at least 9\% absolute accuracy over all available open-source models as well as most closed-source alternatives (such as Inflection-2 and Gemini Pro). This advantage is maintained even compared to much larger models (for instance, Qwen 72B) or those that have been extensively tuned via math-specific reinforcement learning methods (e.g., WizardMath-v1.1 7B). Although DeepSeekMath-Instruct’s performance on MATH is comparable to leading Chinese proprietary models like GLM-4 and Baichuan-3, it remains behind models such as GPT-4 and Gemini Ultra.

When models are evaluated in a setting that permits both natural language reasoning and the use of programmatic tools to tackle problems, DeepSeekMath-Instruct 7B nearly achieves a 60\% accuracy rate on MATH, outperforming every previously available open-source model. On additional benchmarks, the performance of our model rivals that of DeepSeek-LLM-Chat 67B, the former state-of-the-art, despite being only one-tenth its size.

\subsection{Group Relative Policy Optimization} Following the Supervised Fine-Tuning (SFT) phase, reinforcement learning (RL) has shown success in advancing the mathematical reasoning capabilities of LLMs \citep{wang2023math,wizardmath}. Here, we present Group Relative Policy Optimization (GRPO), our novel RL approach designed to be both efficient and effective.

\subsubsection{From PPO to GRPO} Proximal Policy Optimization (PPO) \citep{schulman2017proximal} is an actor-critic RL algorithm that is widely used in the RL fine-tuning stage of LLMs \citep{ouyang2022training}. In particular, it optimizes LLMs by maximizing the following surrogate objective:

\begin{equation}
\footnotesize
    \mathcal{J}_{PPO}(\theta) = \mathbb{E}{[q \sim P(Q),\ o \sim \pi_{\theta_{old}}(O|q)]} \frac{1}{|o|} \sum_{t=1}^{|o|} \min \left\{ \frac{\pi_\theta(o_{t} | q, o_{<t})}{\pi_{\theta_{old}}(o_{t} | q, o_{<t})} A_{t}, \text{clip} \left( \frac{\pi_\theta(o_{t} | q, o_{<t})}{\pi_{\theta_{old}}(o_{t} | q, o_{<t})}, 1 - \varepsilon, 1 + \varepsilon \right)  A_{t} \right\} ,
\end{equation}

The Proximal Policy Optimization (PPO) objective function, denoted as $\mathcal{J}_{PPO}(\theta)$, is formally defined above. The expectation is taken over queries $q$ sampled from $P(Q)$ and outputs $o$ generated by the prior policy $\pi_{\theta_{old}}(O|q)$. For each output sequence, the mean is computed over all time steps $t$. At each position, the surrogate loss is given by the minimum of two terms: the product of the policy ratio $\frac{\pi_\theta(o_{t} | q, o_{<t})}{\pi_{\theta_{old}}(o_{t} | q, o_{<t})}$ and the advantage estimate $A_{t}$, and the product of the advantage $A_{t}$ with the same policy ratio that has been clipped within the interval $[1-\varepsilon, 1+\varepsilon]$.

In this context, $\pi_{\theta}$ denotes the updated policy model, while $\pi_{\theta_{old}}$ corresponds to the previous version. The variables $q$ and $o$ refer to questions and outputs, which are drawn from the dataset of questions and from the old policy $\pi_{\theta_{old}}$, respectively. The parameter $\varepsilon$ serves as a hyperparameter related to clipping in PPO, helping to maintain the stability of the optimization process. $A_t$ represents the advantage term, estimated via Generalized Advantage Estimation (GAE) \citep{gae}, which utilizes the subsequent rewards $\{r_{\ge t}\}$ in conjunction with a learned value function $V_{\psi}$. Consequently, PPO requires the concurrent learning of a value function alongside the policy model. To counteract potential over-optimization with respect to the reward model, a commonly adopted technique is to include a KL-divergence penalty, calculated per token, between the policy and a reference model in the reward signal at every token position \citep{ouyang2022training}, that is,

\begin{equation}
    r_{t} = r_\varphi(q, o_{\le t}) - \beta \log\frac{\pi_{\theta}(o_{t}|q, o_{<t})}{\pi_{ref}(o_{t}|q, o_{<t})},
\label{eq:PPO-reward}
\end{equation}

Here, $r_\varphi$ denotes the reward model, $\pi_{ref}$ represents the reference model—typically the starting SFT model—and $\beta$ stands for the KL penalty coefficient.

As the value function employed in PPO is typically another model of comparable size as the policy model, it brings a substantial memory and computational burden. Additionally, during RL training, the value function is treated as a baseline in the calculation of the advantage for variance reduction. While in the LLM context, usually only the last token is assigned a reward score by the reward model, which may complicate the training of a value function that is accurate at each token. To address this, as shown in Figure \ref{fig:grpo}, we propose Group Relative Policy Optimization (GRPO), which obviates the need for additional value function approximation as in PPO, and instead uses the average reward of multiple sampled outputs, produced in response to the same question, as the baseline. More specifically, for each question $q$, GRPO samples a group of outputs $\{o_1, o_2, \cdots, o_G\}$  from the old policy  $\pi_{\theta_{old}}$  and then optimizes the policy model by maximizing the following objective:

\begin{equation}
\footnotesize
\begin{split}
    \mathcal{J}_{GRPO}(\theta) &= \mathbb{E}{[q \sim P(Q), \{o_i\}_{i=1}^G \sim \pi_{\theta_{old}}(O|q)]}  \\
    & \frac{1}{G}\sum_{i=1}^G\frac{1}{|o_i|} \sum_{t=1}^{|o_i|} \left\{ \min \left[ \frac{\pi_\theta(o_{i,t} | q, o_{i,<t})}{\pi_{\theta_{old}}(o_{i,t} | q, o_{i,<t})} \hat{A}_{i,t}, \text{clip} \left( \frac{\pi_\theta(o_{i,t} | q, o_{i,<t})}{\pi_{\theta_{old}}(o_{i,t} | q, o_{i,<t})}, 1 - \varepsilon, 1 + \varepsilon \right)  \hat{A}_{i,t} \right] - \beta \mathbb{D}_{KL}\left[\pi_{\theta} || \pi_{ref}\right]\right\} ,
\end{split}
\label{eq:GRPO-obj}
\end{equation}

where $\varepsilon$ and $\beta$ are hyper-parameters, and $\hat{A}_{i,t}$  is the advantage calculated based on relative rewards of the outputs inside each group only, which will be detailed in the following subsections. The group relative way that GRPO leverages to calculate the advantages, aligns well with the comparative nature of rewards models,  as reward models are typically trained on datasets of comparisons between outputs on the same question. Also note that, instead of adding KL penalty in the reward, GRPO regularizes by directly adding the KL divergence between the trained policy and the reference policy to the loss, avoiding complicating the calculation of  $\hat{A}_{i,t}$. And different from the KL penalty term used in (\ref{eq:PPO-reward}), we estimate the KL divergence with the following unbiased estimator \citep{kl_approx}:

\begin{equation}
\small
    \mathbb{D}_{KL}\left[\pi_{\theta} || \pi_{ref}\right] = \frac{\pi_{ref}(o_{i,t}|q,o_{i,<t})}{\pi_{\theta}(o_{i,t}|q,o_{i,<t})} - \log\frac{\pi_{ref}(o_{i,t}|q,o_{i,<t})}{\pi_{\theta}(o_{i,t}|q,o_{i,<t})} - 1,
\end{equation}

which is guaranteed to be positive.

\begin{algorithm}[t]
  \small
  \caption{Iterative Group Relative Policy Optimization}
  \textbf{Input} initial policy model $\pi_{\theta_{\text{init}}}$; reward models $r_\varphi$; task prompts $\mathcal{D}$; 
  hyperparameters $\varepsilon$, $\beta$, $\mu$
  \begin{algorithmic}[1]
    \State Initialize policy model $\pi_\theta = \pi_{\theta_{\text{init}}}$
    \For{iteration = 1, \dots, I}
       \State Set the reference model as $\pi_{ref} = \pi_{\theta}$
      \For{step = 1, \dots, M}
      \State Draw a batch $\mathcal{D}_b$ from $\mathcal{D}$
      \State Assign the current policy model to $\pi_{\theta_{old}} = \pi_{\theta}$
      \State For every question $q \in \mathcal{D}_b$, generate $G$ outputs $\{o_i\}_{i=1}^G$ by sampling from $\pi_{\theta_{old}} (\cdot \mid q) $
      \State Evaluate each sampled output $o_i$ to compute its corresponding reward $r_i$ using the reward model $r_{\varphi}$
      \State For each sampled output, determine group relative advantage $\hat{A}_{i,t}$ for token $t$ within $o_i$
      \For{GRPO iteration = 1, \dots, $\mu$}
        \State Perform policy update on $\pi_{\theta}$ by maximizing the group-relative policy objective as specified in Equation \ref{eq:GC-GRPO}
      \EndFor
    \EndFor
    \State Continually refine $r_\varphi$ through replay-based training.
    \EndFor 
  \end{algorithmic}
  \textbf{Output} $\pi_\theta$
  \label{alg:iter-grpo}
\end{algorithm}

\subsubsection{Outcome Supervision RL with GRPO} In this approach, given a question $q$, a set of candidate outputs $\{o_1, o_2, \cdots, o_G\}$ is generated from the previous policy model $\pi_{\theta_{old}}$. Each generated output is assigned a reward through a reward model, resulting in a collection of $G$ rewards, denoted as $\mathbf{r}=\{r_1, r_2, \cdots, r_G\}$. To standardize these rewards, the mean value of the group is subtracted from each reward, followed by division by the standard deviation within the group. Under outcome supervision, the normalized reward is attributed to the terminal token of every output $o_i$, and this normalized value is uniformly used as the advantage $\hat{A}_{i, t}$ for all tokens comprising that output, specifically, $\hat{A}_{i, t} = \widetilde{r}_i = \frac{r_i- {\rm mean}(\mathbf{r})}{{\rm std}(\mathbf{r})}$. The policy is subsequently refined by maximizing the objective stated in equation (\ref{eq:GRPO-obj}).

\subsubsection{Process Supervision RL with GRPO} 

Relying solely on outcome supervision, which assigns a reward only at the conclusion of an output, often proves inadequate and inefficient for guiding the policy on complex mathematical problems. In alignment with \cite{wang2023math}, our work investigates process supervision, where feedback is provided at each individual reasoning step. Formally, for a given question $q$ and a set of $G$ sampled outputs $\{o_1, o_2, \cdots, o_G\}$, a process reward model evaluates each step and produces corresponding rewards, resulting in: $\mathbf{R} = \{ \{r_1^{index(1)},\cdots,r_1^{index(K_1)}\}, \cdots,  \{r_G^{index(1)},\cdots,r_G^{index(K_G)}\} \}$, where $index(j)$ denotes the end token index for the $j$-th step and $K_i$ is the number of reasoning steps in the $i$-th output. To standardize the scale of these rewards, we perform normalization using their mean and standard deviation: $\widetilde{r}_i^{index(j)} = \frac{r_i^{index(j)} - {\rm mean(\mathbf{R})}}{{\rm std(\mathbf{R})}}$.

The process supervision framework then determines the advantage at each token position by aggregating the normalized rewards associated with all subsequent steps, specifically $\hat{A}_{i, t} = \sum_{index(j) \ge t} \widetilde{r}_i^{index(j)}$. Finally, policy optimization is carried out by maximizing the objective outlined in equation (\ref{eq:GRPO-obj}).

\subsubsection{Iterative RL with GRPO} During the course of reinforcement learning, the previous reward model may no longer provide effective guidance for the evolving policy model. To address this, we investigate an iterative RL framework utilizing GRPO. As detailed in Algorithm \ref{alg:iter-grpo}, this approach involves producing updated training data for the reward model using samples drawn from the current policy model. The reward model is further optimized with a replay strategy that maintains 10\% of past data within each training batch. Subsequently, the reference model is updated to be the current policy model, which is then further trained in conjunction with the newly improved reward model.

\subsection{Training and Evaluating DeepSeekMath-RL}

Our reinforcement learning process is centered on DeepSeekMath-Instruct 7B. For RL training, we employ questions in the chain-of-thought format from the SFT dataset that pertain to GSM8K and MATH, totaling approximately 144,000 examples. We deliberately omit all other SFT data in order to analyze the effects of RL on benchmarks that remain data-scarce throughout RL training. The dataset for the reward models is prepared in accordance with \citep{wang2023math}. The initial reward model is trained using DeepSeekMath-Base 7B with a learning rate of $2 \times 10^{-5}$. When applying GRPO, the learning rate for the policy model is set at $1 \times 10^{-6}$, and a KL divergence coefficient of 0.04 is used. For every question, $64$ samples are generated. The maximum sequence length is restricted to 1024, with a batch size of 1024 during training. Following each exploration phase, only one update is performed on the policy model. In evaluation, we benchmark DeepSeekMath-RL 7B using the same criteria as DeepSeekMath-Instruct 7B. Under this setup, tasks involving chain-of-thought reasoning for GSM8K and MATH are considered in-domain for DeepSeekMath-RL 7B, while all other test sets are classified as out-of-domain.

Table \ref{tab:sft_rl_math} presents comparative results for both open- and closed-source models employing chain-of-thought and tool-integrated reasoning approaches across English and Chinese benchmarks. Our analysis reveals: 1) Utilizing chain-of-thought reasoning, DeepSeekMath-RL 7B achieves accuracy rates of 88.2\% on GSM8K and 51.7\% on MATH. These scores not only exceed those of all existing open-source models within the 7B to 70B parameter range, but also outperform most closed-source counterparts. 2) Notably, DeepSeekMath-RL 7B’s training was limited exclusively to chain-of-thought formatted instruction tuning data for GSM8K and MATH, building upon DeepSeekMath-Instruct 7B. In spite of this narrowly focused training set, DeepSeekMath-RL 7B demonstrates superior results compared to DeepSeekMath-Instruct 7B in every evaluation measure, underscoring the advantages conferred by reinforcement learning.

\section{Discussion}
Here, we present the insights gained from our investigations during both the pre-training phase and the reinforcement learning (RL) experiments.

\subsection{Lessons Learnt in Pre-Training}

To begin, we discuss our observations from the pre-training phase. Except where stated otherwise, our training configurations follow those detailed in Section \ref{sec:quality-policy}. Additionally, for the purposes of this section, references to the DeepSeekMath Corpus pertain to a dataset comprising 89 billion tokens, obtained from the second cycle of data gathering.

\subsubsection{Code Training Benefits Mathematical Reasoning}

A widely discussed, though unproven, hypothesis proposes that training on code enhances reasoning skills. In this work, we seek to partially address this idea, focusing on mathematics: training with code boosts models' mathematical reasoning capabilities, whether or not they utilize external tools.

To study how code training affects mathematical reasoning, we experimented with the following two-stage training and one-stage training settings:

\textbf{Two-Stage Training}

\begin{itemize}[topsep=0pt]
    \item \textbf{Code Training for 400B Tokens $\rightarrow$ Math Training for 150B Tokens}: We initially expose DeepSeek-LLM 1.3B to 400B tokens of code, and then subsequently train it on 150B tokens of mathematical data.
    \item \textbf{General Training for 400B Tokens $\rightarrow$ Math Training for 150B Tokens}: As a comparative baseline, we substitute the code data in the initial stage with general-purpose tokens—drawn from an extensive general dataset compiled by DeepSeek-AI—in order to systematically assess whether pre-training on code yields superior mathematical reasoning capability compared to general text.
\end{itemize}

\textbf{One-Stage Training}

\begin{itemize}[topsep=0pt]
    \item \textbf{Math Pre-training on 150B Tokens}: We pre-train DeepSeek-LLM 1.3B using a dataset comprising 150 billion mathematical tokens;
    \item \textbf{Joint Training on 400B Code Tokens and 150B Math Tokens}: Sequential math training after initial code training leads to diminished performance on code-related tasks. To address this, we explore whether simultaneous exposure to both code and math tokens during a single training phase can enhance mathematical reasoning skills while also reducing catastrophic forgetting in coding abilities.
\end{itemize}

\paragraph{Results}
The downstream outcomes for various training configurations are presented in Table \ref{tab:code-math} and Table \ref{tab:code-math-general-eval}.

Code training supports program-aided mathematical reasoning in both two-stage and one-stage training scenarios. According to Table \ref{tab:code-math}, applying code training exclusively in the two-stage framework already leads to notable performance gains when addressing GSM8K and MATH tasks through Python. Introducing math-specific training in the subsequent stage leads to further advancements. Notably, in the one-stage training context, integrating both code tokens and mathematical tokens during training not only helps alleviate the catastrophic forgetting observed in the two-stage approach but also enhances both coding proficiency (Table \ref{tab:code-math-general-eval}) and performance in program-aided mathematical problem solving (Table \ref{tab:code-math}).

Code training independently advances mathematical reasoning abilities even in the absence of tool use. In a two-stage training framework, commencing with code training alone already yields noticeable improvements. This approach further increases the effectiveness of subsequent mathematical training, ultimately achieving optimal results. Conversely, attempting to merge code tokens with math tokens in a single-stage training regimen diminishes mathematical reasoning when tools are not employed. A possible explanation is that the relatively small scale of DeepSeek-LLM 1.3B restricts its capacity to effectively integrate code and mathematical information when exposed to both concurrently.

\subsubsection{ArXiv Papers Seem Ineffective in Improving Mathematical Reasoning}

ArXiv papers are commonly included as a component of math pre-training data \citep{minerva,gpt-f,llemma,mathpile}. However, detailed analysis regarding their impact on mathematical reasoning has not been extensively conducted. Perhaps counter-intuitively, according to our experiments, arXiv papers seem ineffective in improving mathematical reasoning. We experiment with models of different sizes, including DeepSeek-LLM 1.3B and DeepSeek-Coder-Base-v1.5 7B \citep{deepseek-coder}, using arXiv corpora that underwent varied processing pipelines:

\begin{itemize}[topsep=0pt]
    \item \textbf{MathPile} \citep{mathpile}: a dataset comprising 8.9 billion tokens assembled using specific heuristics for data cleaning and selection, with scientific arXiv papers representing more than 85\% of its content;
    \item \textbf{ArXiv-RedPajama} \citep{redpajama}: a collection containing all arXiv LaTeX documents, with the preambles, commentaries, macro definitions, and bibliographic sections excised, resulting in a dataset of 28.0 billion tokens.
\end{itemize}

During our experimentation, we independently trained DeepSeek-LLM 1.3B for 150B tokens and DeepSeek-Coder-Base-v1.5 7B for 40B tokens using each respective arXiv dataset. Our findings indicate that arXiv papers contribute minimally to enhancements in mathematical reasoning abilities. Models trained exclusively on arXiv datasets exhibit either stagnation or a decline in performance across a range of mathematical evaluation tasks of varying difficulty featured in this research. These evaluations encompass quantitative reasoning datasets such as GSM8K and MATH (see Table \ref{tab:arxiv-cot}), multiple-choice assessments like MMLU-STEM (see Table \ref{tab:arxiv-cot}), as well as formal mathematics tasks like miniF2F (see Table \ref{tab:arxiv_atp}).

However, this conclusion has its limitations and should be taken with a grain of salt. We have not yet studied:

\begin{itemize}[topsep=0pt]
    \item How arXiv tokens influence particular mathematical tasks outside the scope of this work, for example, the transformation of formal theorems or proofs into their informal counterparts;
    \item The influence of arXiv tokens when utilized alongside additional datasets;
    \item If the advantages observed from arXiv documents would also appear when deploying models of greater scale.
\end{itemize}

Thus, further exploration is required, which we leave for future studies.

\subsection{Insights of Reinforcement Learning}
\subsubsection{Towards to a Unified Paradigm} In this part, we introduce a cohesive framework for examining a variety of training approaches, including SFT, RFT, DPO, PPO, and GRPO. Additionally, we perform empirical studies to investigate the components influencing this unified approach. In broad terms, the gradient of a training objective with respect to the parameter $\theta$ can be expressed as follows:

\begin{equation}
    \nabla_{\theta}\mathcal{J}_{\textcolor{red}{\mathcal{A}}}(\theta) = \mathbb{E}[\underbrace{(q,o) \sim \textcolor{red}{\mathcal{D}}}_{Data \ Source}]\left( \frac{1}{|o|} \sum_{t=1}^{|o|}  \underbrace{GC_{{\mathcal{A}}}(q, o, t, \textcolor{red}{\pi_{{rf}}})}_{Gradient \ Coefficient}  \nabla_{\theta}\log \pi_{\theta}(o_t | q, o_{<t})\right).
\label{eq:objective}
\end{equation}

There exist three key components: 1) \textit{Data Source $\mathcal{D}$}, which determines the training data; 2) \textit{Reward Function $\pi_{{rf}}$}, which is the source of the training reward signal; 3) \textit{Algorithm $\mathcal{A}$}: which processes the training data and the reward signal to the gradient coefficient $GC$ that determines the magnitude of the penalty or reinforcement for the data. We analyze several representative methods based on such a unified paradigm:

\begin{itemize}
    \item  \textbf{Supervised Fine-tuning (SFT)}: In SFT, a pretrained model undergoes additional training using curated datasets that have been selected by humans.
    \item \textbf{Rejection Sampling Fine-tuning (RFT)}: RFT involves further training of the SFT model by utilizing a subset of outputs generated from the SFT model itself, in response to SFT prompts. Only outputs that produce correct answers are retained and used for further fine-tuning.
    \item \textbf{Direct Preference Optimization (DPO)}: DPO enhances the SFT model by conducting further fine-tuning on augmented data, composed of paired outputs sampled from the SFT model. Optimization is performed using a pairwise DPO loss function.
    \item \textbf{Online Rejection Sampling Fine-tuning (Online RFT)}: In contrast to standard RFT, Online RFT begins with the SFT model as the initial policy, and then continually updates it by fine-tuning with new augmented outputs generated in real time from the current policy model.
    \item \textbf{PPO/GRPO}: With PPO/GRPO, the process starts with initializing the policy model from the SFT model, which is then further improved through reinforcement learning using outputs sampled in real time from the current policy model.
\end{itemize}

Table \ref{tab:data_policy} presents an overview of the constituent elements of these approaches. For an expanded explanation of the derivation procedure, consult Appendix \ref{app:analysis-rl}.

\paragraph{Observation about Data Source} The data source can be categorized into two types: online sampling and offline sampling. Online sampling refers to using training data that originates from the exploratory outputs generated by the current policy model during training. In contrast, offline sampling utilizes training data acquired from the outputs of the initial SFT model. Methods such as RFT and DPO employ offline sampling, whereas Online RFT and GRPO utilize online sampling.

As depicted in Figure \ref{fig:rl-analysis}, our results indicate that Online RFT achieves notably better performance than RFT across both benchmarks. Initially, Online RFT and RFT perform at similar levels; however, as training progresses, Online RFT establishes a clear lead, illustrating the benefits of online training. This observation aligns with expectations: at the outset, the actor's outputs are very similar to those of the SFT model, leading to only subtle variations in the collected data. As training advances, the disparity between samples generated by the actor increases, making real-time data collection progressively more advantageous.

\paragraph{Observation about Gradient Coefficient} The algorithm updates model parameters by channeling the reward signal into the gradient coefficient. Within our experimental setup, the reward function is categorized as either `Rule' or `Model'. The `Rule' type assesses response quality strictly on the basis of answer correctness, whereas the `Model' approach involves developing a reward model that assigns scores to each response. The reward model itself is trained using data derived from rule-based judgments. As shown in Equations \ref{eq:GC-RFT} and \ref{eq:GC-GRPO}, there is a fundamental distinction between GRPO and Online RFT: GRPO distinctively adapts its gradient coefficient in direct response to the reward values assigned by the reward model, which enables a spectrum of reinforcement and penalization depending on the reward magnitude. In comparison, Online RFT does not possess this capacity; it refrains from penalizing incorrect answers, instead providing identical positive reinforcement to all correct responses irrespective of their merit.

As shown in Figure \ref{fig:rl-analysis}, GRPO outperforms online RFT, underscoring the effectiveness of adjusting positive and negative gradient coefficients. Moreover, GRPO+PS achieves better results than GRPO+OS, which demonstrates the advantages of employing more nuanced, step-sensitive gradient adjustments. Additionally, we investigate the impact of iterative RL by conducting two consecutive iterations in our experiments. Figure \ref{fig:iter-rl} illustrates that iterative RL yields notable performance gains, particularly following the initial iteration.

\subsubsection{Why RL Works?} In this study, we apply reinforcement learning to a selected portion of the instruction tuning dataset, observing notable performance gains over the baseline instruction tuning model. To shed light on the underlying reasons for reinforcement learning's effectiveness, we analyze both Pass@K and Maj@K accuracies for the Instruct and RL models across two distinct benchmarks. As depicted in Figure \ref{fig:maj-pass}, reinforcement learning noticeably boosts Maj@K accuracy, while Pass@K remains largely unchanged. These results suggest that reinforcement learning's primary contribution lies in making the output distribution more reliable; specifically, \textbf{the improvements stem from elevating the likelihood of the correct response within the TopK selections, rather than from strengthening the model's core capabilities.} In a similar vein, \citep{wang2023making} highlight a \textbf{misalignment problem} in reasoning tasks handled by SFT models and demonstrate that implementing various preference alignment methods can enhance reasoning abilities in these models \citep{yuan2023rrhf,song2023preference,wang2023making}.

\subsubsection{How to Achieve More Effective RL?}
\label{sec:effec_rl}
Our findings indicate that RL proves to be highly effective for mathematical reasoning problems. Furthermore, we establish a unified framework that facilitates comprehension of various notable training strategies, interpreting each as either a direct application or a streamlined variant of RL. As outlined in Equation \ref{eq:objective}, the framework involves three fundamental elements: Data Source, Algorithm, and Reward Function. We discuss several prospective avenues for advancing each of these aspects.

\paragraph{Data Source} The origin of data underpins all forms of training methodology. In the realm of RL, data sources specifically denote unlabeled queries paired with outputs produced via sampling from the policy model. Our approach in this work restricts itself to questions drawn exclusively from the instruction tuning phase, utilizing straightforward nucleus sampling for generating responses. We conjecture that this design choice may partly explain why our RL framework yields improvements solely on Maj@K metrics. Going forward, we aim to extend our RL pipeline to accommodate out-of-distribution question prompts, complemented by \textbf{more sophisticated sampling (decoding) strategies} such as tree-search-based techniques \citep{yao2023tree}. Additionally, \textbf{efficient inference methods} \citep{xia-etal-2023-speculative,leviathan2023fast,kwon2023efficient,xia2024unlocking}, which critically influence how effectively policy models explore, are likely to be pivotal as well.

\paragraph{Algorithms} Algorithms utilize both the data and the feedback provided by the reward signal to calculate gradient coefficients, thereby facilitating updates to the model parameters. According to Equation \ref{eq:objective}, existing approaches largely \textbf{TRUST} the reward function to modulate the conditional probability of generating specific tokens. Nevertheless, it is not feasible to guarantee the reward signal's accuracy in all scenarios, particularly in highly complex tasks. For instance, the PRM800K datasets \citep{lightman2023let}, which were labeled with considerable care by experienced annotators, still exhibit around 20\% erroneous annotations\footnote{\url{https://github.com/openai/prm800k/issues/12\#issuecomment-1728491852}}. In light of these challenges, our focus will shift toward reinforcement learning algorithms that maintain robustness despite the presence of noisy reward signals. We anticipate that \textbf{WEAK-TO-STRONG} alignment strategies \citep{burns2023weak} could catalyze a significant paradigm shift in learning algorithm design.

\paragraph{Reward Function} The reward function serves as the primary source for generating the training signal. Within the context of RL, the reward is typically modeled by a neural-based reward model. We identify three pivotal challenges associated with reward models: 1) \textbf{Enhancing the generalization capacity of reward models.} A robust reward model must generalize well to unseen, out-of-distribution queries and to novel or more sophisticated decoding results. Failure to do so could result in reinforcement learning that only maintains the status quo of LLM behavior, without delivering true improvements in model capabilities; 2) \textbf{Capturing uncertainty within reward models.} Proper modeling of uncertainty can be instrumental in connecting suboptimal reward models with algorithms that transition from weaker to stronger performance; 3) \textbf{Constructing efficient, high-quality process reward models} that supply detailed, fine-grained training signals for step-wise reasoning procedures \citep{lightman2023let,wang2023math}.

\section{Conclusion, Limitation, and Future Work}  
In this work, we introduce DeepSeekMath, a model that surpasses all open-source models on the MATH benchmark at the competition level, achieving near-parity with proprietary models. DeepSeekMath~is built upon DeepSeek-Coder-v1.5 7B and receives an additional 500B tokens of continued pretraining, with 120B of those tokens consisting of mathematical content extracted from Common Crawl. Through comprehensive ablation studies, we demonstrate that web-sourced pages are a rich source of high-quality mathematical data, whereas arXiv data contribute less than anticipated. We propose Group Relative Policy Optimization (GRPO), an adaptation of Proximal Policy Optimization (PPO), enabling improved mathematical reasoning while lowering memory requirements. Our experimental evaluation confirms the efficacy of GRPO, showing its benefit even when DeepSeekMath-Instruct 7B already achieves strong benchmark results. Furthermore, we offer a unified perspective to interpret various existing approaches and identify several promising research avenues for advancing reinforcement learning techniques.

While DeepSeekMath~performs admirably on tasks involving quantitative reasoning, its effectiveness diminishes in areas such as geometry and theorem-proving when compared with closed-source models. For example, our preliminary experiments reveal that the model struggles with questions pertaining to triangles and ellipses, suggesting a possible bias in the selection of data during both pre-training and fine-tuning stages. Moreover, due to its limited model size, DeepSeekMath~does not match GPT-4 in terms of few-shot learning capability. Whereas GPT-4 demonstrates notable performance gains with few-shot inputs, DeepSeekMath~shows little distinction between zero-shot and few-shot settings. Moving forward, we intend to enhance our engineered data selection strategy to assemble a more robust and higher-quality pre-training dataset. Additionally, we aim to investigate novel methods (Section \ref{sec:effec_rl}) for reinforcing LLMs through more effective reinforcement learning techniques.

\end{document}